% T4 candidate (sonnet3) for fig:barrier — same two-panel Eyring landscape as the existing
% figure, but the connecting curve is now an explicit, disclosed raised-cosine interpolation
% between the exact analytic well/barrier heights of eq:bv (T4_calc.py), not a freehand
% "plot smooth" guess; chi is drawn as an actual horizontal split of the reaction coordinate
% (chi : 1-chi), and a side note ties the same chi-partition to Delta H_a^eff (eq:dHeff).
\begin{tikzpicture}[x=2.6cm,y=1.6cm]
% ================= (a) equilibrium: A=0, chi=1/2 =================
\draw[-{Latex[length=1.4mm]}] (-1.35,-0.15) -- (1.35,-0.15) node[below,font=\scriptsize] {reaction coordinate};
\draw[-{Latex[length=1.4mm]}] (-1.35,-0.15) -- (-1.35,1.25) node[above,font=\scriptsize] {free energy};
\draw[thick] plot[smooth] coordinates {(-1.0000,0.0000) (-0.9167,0.0170) (-0.8333,0.0670) (-0.7500,0.1464) (-0.6667,0.2500) (-0.5833,0.3706) (-0.5000,0.5000) (-0.4167,0.6294) (-0.3333,0.7500) (-0.2500,0.8536) (-0.1667,0.9330) (-0.0833,0.9830) (0.0000,1.0000) (0.0833,0.9830) (0.1667,0.9330) (0.2500,0.8536) (0.3333,0.7500) (0.4167,0.6294) (0.5000,0.5000) (0.5833,0.3706) (0.6667,0.2500) (0.7500,0.1464) (0.8333,0.0670) (0.9167,0.0170) (1.0000,0.0000)};
\fill (-1,0) circle (0.6pt);
\fill (1,0) circle (0.6pt);
\fill (0,1.0) circle (0.6pt);
\node[below,font=\scriptsize] at (-1,-0.02) {filled site};
\node[below,font=\scriptsize] at (1,-0.02) {empty site $+$ ion};
\node[above,font=\scriptsize] at (0,1.02) {transition state};
\draw[densely dotted,gray] (-1,0) -- (-0.55,0);
\draw[densely dotted,gray] (0,1.0) -- (-0.55,1.0);
\draw[{Bar[width=1.6mm]}-{Bar[width=1.6mm]}] (-0.55,0.02) -- (-0.55,0.98)
  node[pos=0.5,left=1pt,font=\scriptsize] {$\Delta G_a$};
\node[font=\scriptsize,align=center] at (0,-0.42) {(a) equilibrium\\$A=0$, $\chi=\tfrac12$};
% ================= (b) driven: A=0.5(scale), chi=0.35 (early TS) =================
\begin{scope}[xshift=7.9cm]
\draw[-{Latex[length=1.4mm]}] (-1.35,-0.15) -- (1.35,-0.15) node[below,font=\scriptsize] {reaction coordinate};
\draw[thick] plot[smooth] coordinates {(-1.0000,0.0000) (-0.9417,0.0141) (-0.8833,0.0553) (-0.8250,0.1208) (-0.7667,0.2062) (-0.7083,0.3057) (-0.6500,0.4125) (-0.5917,0.5193) (-0.5333,0.6187) (-0.4750,0.7042) (-0.4167,0.7697) (-0.3583,0.8109) (-0.3000,0.8250) (-0.1917,0.8024) (-0.0833,0.7362) (0.0250,0.6310) (0.1333,0.4938) (0.2417,0.3340) (0.3500,0.1625) (0.4583,-0.0090) (0.5667,-0.1687) (0.6750,-0.3060) (0.7833,-0.4112) (0.8917,-0.4774) (1.0000,-0.5000)};
\fill (-1,0) circle (0.6pt);
\fill (1,-0.5) circle (0.6pt);
\fill (-0.3,0.825) circle (0.6pt);
\node[below,font=\scriptsize] at (-1,-0.02) {filled site};
\node[below,font=\scriptsize] at (1,-0.52) {empty site $+$ ion};
\node[above,font=\scriptsize] at (-0.3,0.845) {transition state};
% forward barrier Delta G_a - chi A
\draw[densely dotted,gray] (-1,0) -- (-0.62,0);
\draw[densely dotted,gray] (-0.3,0.825) -- (-0.62,0.825);
\draw[{Bar[width=1.6mm]}-{Bar[width=1.6mm]}] (-0.62,0.02) -- (-0.62,0.805)
  node[pos=0.5,left=1pt,font=\scriptsize,align=right] {$\Delta G_a$\\[-1pt]$-\chi A$};
% reverse barrier Delta G_a + (1-chi) A
\draw[densely dotted,gray] (-0.3,0.825) -- (0.15,0.825);
\draw[densely dotted,gray] (1,-0.5) -- (0.15,-0.5);
\draw[{Bar[width=1.6mm]}-{Bar[width=1.6mm]}] (0.15,0.805) -- (0.15,-0.48)
  node[pos=0.5,right=1pt,font=\scriptsize,align=left] {$\Delta G_a$\\[-1pt]$+(1{-}\chi)A$};
% total driving force A (well-to-well drop)
\draw[densely dotted,gray] (1,0) -- (1.18,0);
\draw[densely dotted,gray] (1,-0.5) -- (1.18,-0.5);
\draw[{Bar[width=1.4mm]}-{Bar[width=1.4mm]}] (1.18,-0.02) -- (1.18,-0.48)
  node[pos=0.5,right=1pt,font=\scriptsize] {$A$};
% chi-split dimension line under the axis
\draw[{Bar[width=1.2mm]}-{Bar[width=1.2mm]}] (-1,-0.42) -- (-0.3,-0.42)
  node[pos=0.5,below,font=\scriptsize] {$\chi$};
\draw[{Bar[width=1.2mm]}-{Bar[width=1.2mm]}] (-0.3,-0.42) -- (1,-0.42)
  node[pos=0.5,below,font=\scriptsize] {$1-\chi$};
\node[font=\scriptsize,align=center] at (0,-0.66) {(b) driven $A>0$ ($\chi=0.35$, early TS)};
\node[draw,rounded corners=1pt,font=\tiny,align=center,inner sep=2.2pt,fill=black!4] at (0.35,-1.05)
  {same $\chi$ split recurs with $\Omega$:\\$\Delta H_a^\eff=\Delta H_a-\chi_d\Omega$ (eq.~\eqref{eq:dHeff})};
\end{scope}
\end{tikzpicture}
\caption{활성화 장벽 그림 — 연결 곡선은 $\Delta G_a\mp\chi A$ 같은 지정 높이(식~\eqref{eq:bv})를 잇는
raised-cosine 보간(스켈레톤 T4\_calc.py, freehand 곡선 아님). (a) 평형($A{=}0$) — 두 골짜기 같은 높이,
장벽 $\Delta G_a$, 속도는 Eyring 식. (b) 구동력 $A>0$($\chi{=}0.35$, 반응좌표 위 $\chi\!:\!1{-}\chi$ 분할로
표시된 이른 전이상태)가 정방향 장벽을 $\Delta G_a-\chi A$ 로 낮추고 역방향을 $\Delta G_a+(1-\chi)A$ 로
높인다(두 속도의 비가 식~\eqref{eq:db} detailed balance). 우물 낙차가 구동력 $A$ 자체다. 곁상자 — 같은
$\chi$ 분할이 \S\ref{sec:lag} 의 정규용액 항에도 재현되어 $\Delta H_a^\eff=\Delta H_a-\chi_d\Omega$(식~\eqref{eq:dHeff})를
낳는다. 이 그림이 식~\eqref{eq:xieq} logistic 과 \S\ref{sec:lag} 의 $L_q$ 두 결과의 공통 출발점이다.}
\label{fig:barrier}
